\section*{Exact biased-child coherent identity and collision asymptotics}

All logarithms are natural.  Put
\[
 \psi(t)=t^2\log\frac1{t^2},\qquad
 \mathcal E_{\rm raw}(c)=\sum_T\psi(c_T),
\]
with \(\psi(0)=0\), and define the scalar interference kernel
\[
 K(x,y)=\psi(x)+\psi(y)-\frac{\psi(x+y)+\psi(x-y)}2.
\]

Let \(a,b\) be functions on a block \(W\), let \(h\) be a Boolean sign
on a disjoint block \(X\), and suppose \(a+b\) and \(a-b\) are Boolean
signs.  Thus
\[
 F(W,X)=a(W)+h(X)b(W)
\]
is Boolean.  Fix \(i\in X\), write
\[
 h=r+x_i s,\qquad \mu=\mathbb Eh=\widehat r(\varnothing),
 \qquad \alpha=\widehat h(\{i\})=\widehat s(\varnothing),
\]
and put \(q=\|b\|_2^2=\sum_T b_T^2\).  For any Boolean sign
\(f=u+x_i v\), define its physical split decrement by
\[
 D_i(f)=\mathcal E_{\rm raw}(u)+\mathcal E_{\rm raw}(v)
 -\frac{\mathcal E_{\rm raw}(u+v)+\mathcal E_{\rm raw}(u-v)}2.
\]

\subsection*{Theorem 1 (proved): exact biased-child remainder}

For every choice above,
\[
 D_i(F)=qD_i(h)+R_\mu(\alpha),                                      \tag{1}
\]
where
\begin{align*}
 R_\mu(\alpha)={}&\mathcal E_{\rm raw}(a+\mu b)
 +\mathcal E_{\rm raw}(\alpha b)\\
 &-\frac{\mathcal E_{\rm raw}(a+(\mu+\alpha)b)
              +\mathcal E_{\rm raw}(a+(\mu-\alpha)b)}2
 -qK(\mu,\alpha).                                                   \tag{2}
\end{align*}
Equivalently,
\begin{align*}
 R_\mu(\alpha)={}&\mathcal E_{\rm raw}(a+\mu b)
 +\alpha^2\mathcal E_{\rm raw}(b)\\
 &-\frac{\mathcal E_{\rm raw}(a+(\mu+\alpha)b)
              +\mathcal E_{\rm raw}(a+(\mu-\alpha)b)}2\\
 &+q\left[-\psi(\mu)+\frac{\psi(\mu+\alpha)+\psi(\mu-\alpha)}2\right]. \tag{3}
\end{align*}
When \(\mu=0\), (3) is exactly the reviewed balanced formula, including
\(q\alpha^2\log(1/\alpha^2)\) and both coherent spectra.

\paragraph{Proof.}
The physical split of \(F\) is
\[
 F=(a+rb)+x_i(sb).
\]
For every nonempty Fourier index \(U\) of the remaining child block, the
four vectors entering \(D_i(F)\) are respectively
\(r_Ub,s_Ub,(r_U+s_U)b,(r_U-s_U)b\).  Tensor entropy gives
\[
 \mathcal E_{\rm raw}(cb)=q\psi(c)+c^2\mathcal E_{\rm raw}(b).
\]
The coefficients multiplying \(\mathcal E_{\rm raw}(b)\) cancel because
\[
 r_U^2+s_U^2=\frac{(r_U+s_U)^2+(r_U-s_U)^2}{2}.
\]
The resulting contribution is exactly \(q\) times the corresponding
nonconstant-index contribution to \(D_i(h)\).  At the empty child index,
the four vectors are
\[
 a+\mu b,\quad \alpha b,\quad a+(\mu+\alpha)b,
 \quad a+(\mu-\alpha)b.
\]
Subtracting \(q\) times the empty-index contribution
\(K(\mu,\alpha)\) to \(D_i(h)\) proves (1)--(2).  Finally,
\[
 \mathcal E_{\rm raw}(\alpha b)
 =q\psi(\alpha)+\alpha^2\mathcal E_{\rm raw}(b),
\]
which gives (3). \(\square\)

\subsection*{Coefficientwise form}

For \(b_T\ne0\), put \(t_T=a_T/b_T\).  The quadratic coefficient of
\(\log(1/b_T^2)\) cancels inside each kernel, so (2) is exactly
\[
 \boxed{
 R_\mu(\alpha)=\sum_{T:b_T\ne0}b_T^2
 \bigl[K(t_T+\mu,\alpha)-K(\mu,\alpha)\bigr].}                    \tag{4}
\]
This identity isolates the child-mean correction rather than treating a
biased child as balanced.

\subsection*{Theorem 2 (proved): fixed-bias collision law}

Define the exact collision energy
\[
 Z_\mu(a,b)=\sum_{T:\,b_T\ne0,\ a_T+\mu b_T=0}b_T^2.
\]
For fixed finite coefficient arrays \(a,b\):
\[
 \lim_{\alpha\to0}
 \frac{R_\mu(\alpha)}{\alpha^2\log(1/\alpha^2)}
 =-Z_\mu(a,b),\qquad 0<|\mu|<1,                                  \tag{5}
\]
whereas at exact balance
\[
 \lim_{\alpha\to0}
 \frac{R_0(\alpha)}{\alpha^2\log(1/\alpha^2)}
 =q-Z_0(a,b).                                                       \tag{6}
\]

\paragraph{Proof.}
For every fixed \(x\ne0\), Taylor expansion of \(\psi\) away from zero
gives
\[
 K(x,\alpha)=\psi(\alpha)+O_x(\alpha^2),
\]
whereas \(K(0,\alpha)=0\) identically.  If \(\mu\ne0\), the subtracted
kernel \(K(\mu,\alpha)\) contributes one copy of \(\psi(\alpha)\) at
every \(b_T\ne0\).  It cancels when \(t_T+\mu\ne0\), and remains with
negative sign precisely when \(a_T+\mu b_T=0\), proving (5).  If
\(\mu=0\), the subtracted kernel vanishes.  A positive copy remains
precisely when \(a_T\ne0\), proving (6). \(\square\)

Thus a fixed nonzero child bias removes every positive leading
\(\alpha^2\log(1/\alpha^2)\) debt.  Exact coherent collisions produce a
negative logarithmic refund instead.  This statement is not uniform as
\(\mu\to0\) or as a nonzero collision gap tends to zero.  Indeed
\[
 K(c\alpha,\alpha)=\alpha^2K(c,1),
\]
so the balanced positive logarithm reappears when \(\mu\) is comparable
to \(\alpha\).  Uniform control of this two-parameter transition is part
of the remaining hierarchy gap.

\subsection*{Actual Boolean collision control}

Take the two three-bit branches, in lexicographic truth-table order,
\[
 g_+=(-1,-1,-1,-1,-1,-1,-1,1),
\]
\[
 g_-=(-1,-1,-1,1,-1,1,1,-1).
\]
Their half-sum and half-difference Fourier vectors are
\[
 \widehat a=(-1/2,-1/4,-1/4,0,-1/4,0,0,1/4),
\]
\[
 \widehat b=(-1/4,0,0,1/4,0,1/4,1/4,-1/2).
\]
At \(\mu=1/2\), the collision energy is \(Z_{1/2}=1/4\).  For every
\(n\ge4\), choose a Boolean child which is negative on \(2^{n-2}\)
points, with \(2^{n-3}+1\) of them on the \(x_1=+1\) face and
\(2^{n-3}-1\) on the \(x_1=-1\) face.  Then
\[
 \mathbb Eh=\frac12,\qquad \widehat h(\{1\})=-2^{2-n}.
\]
This is a genuine Boolean family, and (5) gives
\[
 \frac{R_{1/2}(\alpha_n)}{\alpha_n^2\log(1/\alpha_n^2)}\longrightarrow-\frac14.
\]
The archived deterministic values through \(n=14\) approach this limit
from below.  The family is a control for a refund, not an FEI
counterexample.

